\section{Base-level alignment}
\label{sec:final-alignment-role}
\label{sec:final-alignment-cpu-contracts}

In this short chapter I want to discuss last parts of the pipeline's algorithms. The first one is the base-level alignment.
Chaining will tell us approximetely where two sequences match, but it isn't sufficient to tell which exact symbols correspond to which. The base-level alignment performs the dynamic programming specified in the introductory sections~\ref{sec:exact-dp} as
now instead of matching entire target, we can just process target subsequence of size similar to the query size.

\info{Conditional CIGAR execution}{
BioseqDB doesn't call base-level alignment if no \path{bioseq.cigar()} is requested~\ref{lst:xdrop-zdrop-strategies}, so
then \path{bioseq.score()} only reflects the chaining score. If \path{bioseq.cigar()} call is present in the SQL query, then score represents results of the final alignment (thus score can differ  depending on the query).
This is explained in the walkthrough appendix~\ref{app:sql-walkthrough}.
}

In BioseqDB I used existing implementation of KSW2 (the algorithm by the author of minimap2 that uses AVX SIMD C intrinsics to implement DP calculations)
from minimap2 and Block Aligner~\cite{minimap2_paper,block_aligner_paper}.
The DP formulations follow what was described already easlier~\ref{sec:exact-dp}, where KSW2 provides various flavours of functions in its API primarily left/right extension and local mode~\ref{fig:alignment-boundary-modes}.
KSW2 uses banded alignment where the maximum distances are fixed (for more details on that see chaining for maximum diagonal distance (bandwith)~\ref{fig:chaining-dotplot})
The block aligner is more complex and uses moving square blocks to adaptively set this size of "bandwith", where we compute scores in the score matrix.

\begin{table}[H]
\centering
\small
\caption{CPU aligners incorporated into BioseqDB.}
\label{tab:final-aligner-comparison}
\setlength{\tabcolsep}{4pt}
\renewcommand{\arraystretch}{1.1}
\begin{tabularx}{\linewidth}{@{}>{\raggedright\arraybackslash}p{.19\linewidth}
  >{\raggedright\arraybackslash}X >{\raggedright\arraybackslash}X@{}}
\toprule
 & \textbf{KSW2} & \textbf{Block Aligner}\\
\midrule
DP matrix where scores are calculated & Fixed diagonal band & Moving square that grows and shrinks repeatedly \\
Gap cost function (single-affine/dual-affine)& \(\min(q_1+e_1\ell,q_2+e_2\ell)\) & \(q_1+e_1\ell\)\\
Special stop conditions (see below) & Z-drop & X-drop after two consecutive low-scoring steps\\
Parameters that affect sizes of searched regions & "bandwith" is the maximum distance from the diagonal for the diagonal band
 & Block size, which is the initial size of the square \\
\bottomrule
\end{tabularx}
\end{table}

I previously introduced X-drop and Z-drop concepts~\ref{lst:xdrop-zdrop-strategies}, but formally they can be formulated as follows.
Let's say \(S^\star\) is the best score reached at \((i^\star,j^\star)\) and we test this best score against the
current score \(S\) at \((i,j)\) (meaning that we are matching $i$th symbol of target vs. $j$th symbol of query):
\begin{equation}
\begin{aligned}
  \text{X-drop:}\quad&S^\star-S>X,\\
  \text{Z-drop:}\quad&S^\star-S>Z+e_{g}\left|(i-i^\star)-(j-j^\star)\right|.
\end{aligned}
\label{eq:zdrop}
\end{equation}
$e_{g}$ here is gap-extension penalty, which in this context means that we can disard all symbols from the best score to the current position by introducing a gap. This means 
that effectively we just stop earlier if the distance from the diagonal gets to large, so practically Z-drop is more restrictive for correct setup.

\clearpage
\section{How Block Aligner follows the alignment}
\label{sec:final-alignment-block}

Block Aligner moves a \(B\times B\) square through the scoring matrix,
and internally compute strips of width \(S=8\) visible on the figure~\ref{fig:block-adaptive-reconstruction} below.
When the algorithm fails to improve best score for number of steps, then it decides to grow the block as seen on image below~\ref{fig:block-adaptive-reconstruction}

\begin{figure}[H]
\centering
\begin{tikzpicture}[x=0.43cm,y=-0.43cm,>=Latex,font=\small]
  \begin{scope}
    \node at (3,-1) {(a) shift right};
    \foreach \k in {0,...,6} {
      \draw[gray!25,very thin] (0,\k) -- (6,\k) (\k,0) -- (\k,6);
    }
    \fill[blue!15] (4,0) rectangle (5,4);
    \draw[gray!70,dashed,thick] (0,0) rectangle (4,4);
    \draw[blue!70!black,very thick] (1,0) rectangle (5,4);
    \draw[->,thick] (1,4.8) -- (2,4.8) node[midway,below] {$S$};
    \node at (3,7) {$(i,j)\mapsto(i,j+S)$};
  \end{scope}
  \begin{scope}[xshift=5cm]
    \node at (3,-1) {(b) grow when best score doesn't improve};
    \foreach \k in {0,...,6} {
      \draw[gray!25,very thin] (0,\k) -- (6,\k) (\k,0) -- (\k,6);
    }
    \fill[teal!15] (0,3) rectangle (3,6);
    \fill[blue!15] (3,0) rectangle (6,6);
    \draw[blue!70!black,very thick] (0,0) rectangle (6,6);
    \draw[gray!70,dashed,thick] (0,0) rectangle (3,3);
    \node at (1.5,4.5) {$\Omega_1$};
    \node at (4.5,3) {$\Omega_2$};
    \node at (3,7) {$B_0\mapsto 2B_0$};
  \end{scope}
  \begin{scope}[xshift=10cm]
    \node at (3,-1) {(c) then shrink to keep bottom-right quarter};
    \foreach \k in {0,...,6} {
      \draw[gray!25,very thin] (0,\k) -- (6,\k) (\k,0) -- (\k,6);
    }
    \fill[blue!15] (3,3) rectangle (6,6);
    \draw[gray!70,dashed,thick] (0,0) rectangle (6,6);
    \draw[blue!70!black,very thick] (3,3) rectangle (6,6);
    \draw[->,thick] (1.2,1.2) -- (2.8,2.8);
    \node at (3,7) {$B\mapsto B/2$};
  \end{scope}
\end{tikzpicture}
\caption{Illustration of 3 steps that Block Aligner repeatedly performs. Growth computes the lower
rectangle \(\Omega_1\) before the right rectangle \(\Omega_2\).}
\label{fig:block-adaptive-reconstruction}
\end{figure}

The pseudocode below is directly taken from Block Aligner paper~\cite{block_aligner_paper}. For more details please consult the projects README~\cite{block_aligner_readme}.
I use here original paper notation, where \(q\) is the query and \(p\) the
target sequence aligner to.
\(G^c_{open},G^c_{close},G^r_{open},G_{ext}\) are the gap penalties.
Function \(\mathrm{COMPUTE\_RECT}((r,c),(w,h))\) is used to compute all the scores in the given rectangle given as parameters, using SIMD version of regular DP scoring formulation.

The important thing in the paper is: \([r,r+h)\times[c,c+w)\); \(Y=\lfloor B/S\rfloor-1\) which defines threshold, when block assumes the score cannot improve and decides to grow.
Note that here $S$ was used by paper's authors as the width of the strip, not the actual score computed. 

% Keep the introduction compact before the full-page reproduction.
\vfill\clearpage
\begin{algorithm}[H]
\caption{Pseudocode for the Block Aligner algorithm.}
\label{alg:adaptive-block-alignment}
\footnotesize
% Restore the published end statements locally; other thesis listings use noend.
\algrenewtext{EndIf}{\algorithmicend\ \algorithmicif}
\algrenewtext{EndWhile}{\algorithmicend\ \algorithmicwhile}
\algrenewtext{EndProcedure}{\algorithmicend\ \algorithmicprocedure}
% Wording (including the paper's typos) and numbered blank lines are intentional.
\begin{algorithmic}[1]
  \Procedure{ALIGN}{$q,p,X,B_{min},B_{max},G^c_{open},G^c_{close},G^r_{open},G_{ext}$}
    \State $S\gets8$, $\mathit{max\_score}\gets0$, $\mathit{dir}\gets\textsc{GROW}$,
      $B\gets B_{min}$, $B_{prev}\gets0$, $i\gets0$, $j\gets0$
    \State Initialize $D$, $D'$, $C$, $R$, and the trace array
    \While{true}
      \Comment{\parbox[t]{.69\linewidth}{\raggedright One step of ligner. Note that $i$, $j$, $B$, and $\mathit{dir}$ have already been updated.}}
      \If{$\mathit{dir}=\textsc{RIGHT}$}
        \Comment{Block position $(i,j-S)\mapsto(i,j)$}
        \State $\mathit{curr\_score}\gets\Call{COMPUTE\_RECT}{(i,j+B-S),(S,B)}$
      \ElsIf{$\mathit{dir}=\textsc{DOWN}$}
        \Comment{Block position $(i-S,j)\mapsto(i,j)$}
        \State $\mathit{curr\_score}\gets\Call{COMPUTE\_RECT}{(i+B-S,j),(B,S)}$
      \ElsIf{$\mathit{dir}=\textsc{GROW}$}
        \State $\mathit{curr\_score}\gets\Call{COMPUTE\_RECT}{(i+B_{prev},j),(B_{prev},B-B_{prev})}$
          \Comment{Grow down}
        \State $\mathit{curr\_score}\gets\Call{COMPUTE\_RECT}{(i,j+B_{prev}),(B-B_{prev},B)}$
          \Comment{Grow right}
      \EndIf
      \State \strut
      \If{$\mathit{curr\_score}>\mathit{max\_score}$}
        \State Update $\mathit{max\_score}$ with $\mathit{curr\_score}$ and its store its position
        \State Save $i$, $j$, and the block right and bottom border scores as checkpoint
      \EndIf
      \State \strut
      \State \textbf{if} $\mathit{curr\_score}<\mathit{max\_score}-X$ for 2 consecutive steps, \textbf{then break}
        \Comment{X-drop}
      \State \textbf{if} $i+B>|q|$ \textbf{and} $j+B>|p|$, \textbf{then break}
        \Comment{Reached end of DP matrix}
      \State \textbf{if} $i+B>|q|$ \textbf{or} $j+B>|p|$, \textbf{then} shift $i$, $j$ by $S$ to avoid out of bounds \textbf{and continue}
      \State \strut
      \If{$B<B_{max}$ and $Y$ steps pass without improving $\mathit{max\_score}$}
        \Comment{Block grow criteria}
        \State $B_{prev}=B$, $B\gets2B$, $\mathit{dir}\gets\textsc{GROW}$
        \State Restore $i$, $j$, and block border scores from previous checkpoint
        \State \textbf{continue}
      \EndIf
      \State \strut
      \If{$B>B_{min}$ and $\mathit{max\_score}$ is in the bottom-right corner region of block}
        \Comment{Block shrink criteria}
        \State $B\gets B/2$, $i\gets i+B$, $j\gets j+B$
        \State Save $i$, $j$, and the block right and bottom border scores as checkpoint
      \EndIf
      \State \strut
      \State $\mathit{dir}\gets\textsc{RIGHT}$ or \textsc{DOWN} based on comparing right and bottom border cells
      \State Increment $i$, $j$ by $S$ based on $\mathit{dir}$
    \EndWhile
    \State \Return the trace array, and \textbf{if} X-drop, \textbf{then} $\mathit{max\_score}$, \textbf{else} $D[|q|][|p|]$
  \EndProcedure
\end{algorithmic}
\end{algorithm}


% Other existing illustrations retained as source material.
% \begin{figure}[htbp]
% \centering
% \begin{tikzpicture}[
%   >=Latex,
%   box/.style={draw,rounded corners,fill=blue!7,align=center,
%               minimum width=3.25cm,minimum height=1.15cm,font=\small},
%   note/.style={align=center,font=\small}
% ]
%   \node[box] (left) at (1.9,0) {left extension\\reversed flank};
%   \node[box] (acc) at (5.6,0) {accumulate nearby\\anchor intervals};
%   \node[box] (fill) at (9.3,0) {mandatory global\\fills};
%   \node[box] (right) at (13.0,0) {right extension\\forward flank};
%   \draw[->,thick] (left.east) to (acc.west);
%   \draw[->,thick] (acc.east) to (fill.west);
%   \draw[->,thick] (fill.east) to (right.west);
%   \node[note] at (7.45,-1.05)
%     {ordered anchors preserve the query and target coordinate direction};
% \end{tikzpicture}
% \caption{Decomposition of a chained region into two extensions and a set of
% global inter-anchor fills. Several short anchor intervals may form one fill.}
% \label{fig:three-phase-alignment}
% \end{figure}

% \begin{figure}[htbp]
% \centering
% \begin{tikzpicture}[x=0.48cm,y=0.48cm,>=Latex,font=\small]
%   \begin{scope}
%     \draw[step=1,gray!24,very thin] (0,0) grid (6,6);
%     \foreach \x in {0,2,4,6} {\draw[blue!65!black,thick] (\x,0) to (\x,6);}
%     \foreach \y in {0,2,4,6} {\draw[blue!65!black,thick] (0,\y) to (6,\y);}
%     \draw[orange!80!black,very thick]
%       (0.1,0.2) .. controls (1.6,1.0) and (1.4,2.0) .. (2.2,2.2)
%       .. controls (3.8,2.5) and (3.5,4.0) .. (4.2,4.1)
%       .. controls (5.0,4.8) and (5.4,5.4) .. (5.9,5.8);
%     \node[align=center,text width=4.8cm] at (3,7.2)
%       {forward pass stores tile borders};
%     \node at (3,-0.9) {(a) boundary state};
%   \end{scope}
%   \begin{scope}[xshift=7.3cm]
%     \draw[step=1,gray!24,very thin] (0,0) grid (6,6);
%     \foreach \x in {0,2,4,6} {\draw[blue!65!black,thick] (\x,0) to (\x,6);}
%     \foreach \y in {0,2,4,6} {\draw[blue!65!black,thick] (0,\y) to (6,\y);}
%     \fill[purple!18] (0,0) rectangle (2,2);
%     \fill[purple!18] (2,2) rectangle (4,4);
%     \fill[purple!18] (4,4) rectangle (6,6);
%     \draw[orange!80!black,very thick,<-]
%       (0.1,0.2) .. controls (1.6,1.0) and (1.4,2.0) .. (2.2,2.2)
%       .. controls (3.8,2.5) and (3.5,4.0) .. (4.2,4.1)
%       .. controls (5.0,4.8) and (5.4,5.4) .. (5.9,5.8);
%     \node[align=center,text width=4.8cm] at (3,7.2)
%       {backward pass recomputes visited tiles};
%     \node at (3,-0.9) {(b) local replay};
%   \end{scope}
% \end{tikzpicture}
% \caption{Gridded traceback. Blue tile borders persist after the forward pass.
% Only the purple tiles crossed by the selected path are recomputed during
% traceback.}
% \label{fig:gridded-traceback}
% \end{figure}
